承载力继续提高后，三级食物链不会一直维持规则振荡，而是会走向更复杂的长期波动。有限时间最大李雅普诺夫指数支持了这一方向，但还不足以给出唯一临界点；文中的 $K$ 也只表示归一化结构压力，而非长江某一实测承载量。污染和入侵则把另一层风险推了出来：食源过程量可能和珍稀物种终态开始分离。按我们的分析框架看，禁渔只是让恢复链条重新启动的起点。链条启动以后，收益能不能留在系统内部，最终还得看水质、连通性、放流和入侵控制能不能跟上。
